\documentclass[a4paper, 12pt]{article}

% --- Pacotes Essenciais ---
\usepackage[utf8]{inputenc}
\usepackage[portuguese]{babel}
\usepackage{geometry}
\usepackage{amsmath} % Pacote principal para matemática
\usepackage{amssymb} % Para símbolos como \mathbb{R}
\usepackage{graphicx}
\usepackage{caption}

% --- Configuração de Página ---
\geometry{a4paper, margin=1in}

% --- Informações do Documento ---
\title{Formulação do Controle em Nível de Posição \\
       para um Manipulador 2-DOF \\
       \large (Baseado em Cinemática Inversa Analítica)}
\author{Gemini}
\date{\today}

% --- Início do Documento ---
\begin{document}

\maketitle

\section{Objetivo e Abordagem}
O objetivo é controlar a posição do efetuador final (EE) de um manipulador 2-DOF para que ele alcance uma posição de referência $\mathbf{x}_d = \begin{bmatrix} x_d \\ y_d \end{bmatrix}$ no espaço cartesiano.

Ao contrário das abordagens de "taxa resolvida" (nível de velocidade) que usam a Jacobiana para calcular velocidades de junta ($\dot{\mathbf{q}}$) iterativamente, esta solução opera no \textbf{nível de posição}.

A estratégia consiste em duas etapas:
\begin{enumerate}
    \item \textbf{Resolver a Cinemática Inversa (CI) Analiticamente:} Encontrar a solução exata $\mathbf{q}_d = \begin{bmatrix} q_{1d} \\ q_{2d} \end{bmatrix}$ que satisfaz a equação de cinemática direta:
    \[
    \mathbf{x}_d = f(\mathbf{q}_d)
    \]
    Esta solução é calculada \textit{uma única vez} quando o alvo é definido.
    \item \textbf{Gerar uma Trajetória (Interpolação):} Criar um caminho suave no espaço de juntas, $\mathbf{q}(t)$, que começa na posição atual $\mathbf{q}_{\text{atual}}$ e termina na posição alvo $\mathbf{q}_d$ ao longo de uma duração $T$.
\end{enumerate}

\section{Etapa 1: Solução da Cinemática Inversa Analítica}
Para um manipulador planar 2-DOF, podemos encontrar uma solução de forma fechada (analítica) usando geometria. A função \texttt{inverse\_kinematics(x, y)} implementa esta solução.

\subsection{Verificação de Alcance}
Primeiro, verificamos se a posição alvo $(x_d, y_d)$ é alcançável. A distância da base ao alvo, $D = \sqrt{x_d^2 + y_d^2}$, deve estar dentro dos limites do robô:
\[
|L_1 - L_2| \le D \le L_1 + L_2
\]
Onde $L_1$ e $L_2$ são os comprimentos dos elos. Se $D$ estiver fora desse intervalo, a função retorna \texttt{None}, pois não existe solução.

\subsection{Cálculo de $q_2$ (Ângulo do Cotovelo)}
Considerando o triângulo formado pela base, o cotovelo e o efetuador final (com lados $L_1$, $L_2$ e $D$), podemos usar a Lei dos Cossenos para encontrar o ângulo $q_2$:
\[
D^2 = L_1^2 + L_2^2 - 2 L_1 L_2 \cos(\pi - |q_2|)
\]
Usando a identidade $\cos(\pi - \theta) = -\cos(\theta)$, e definindo $q_2$ como o ângulo interno, temos:
\[
\cos(|q_2|) = \frac{D^2 - L_1^2 - L_2^2}{2 L_1 L_2}
\]
O código escolhe uma das duas soluções possíveis (cotovelo para cima ou para baixo). A linha \texttt{q2 = -math.acos(...)} seleciona a configuração "cotovelo para cima" (assumindo uma convenção de ângulo).

\subsection{Cálculo de $q_1$ (Ângulo da Base)}
O ângulo $q_1$ é composto por dois ângulos:
\begin{enumerate}
    \item $\alpha$: O ângulo direto da base até o alvo, dado por $\alpha = \text{atan2}(y_d, x_d)$.
    \item $\beta$: O ângulo formado entre o elo $L_1$ e a linha $D$.
\end{enumerate}
Usando a Lei dos Cossenos novamente no mesmo triângulo, podemos encontrar $\beta$:
\[
\cos(\beta) = \frac{D^2 + L_1^2 - L_2^2}{2 L_1 D}
\]
O ângulo final $q_1$ é então a soma ou subtração desses ângulos, dependendo da configuração do cotovelo. Para a solução "cotovelo para cima" escolhida ($q_2$ negativo), temos:
\[
q_1 = \alpha + \beta = \text{atan2}(y_d, x_d) + \arccos\left(\frac{D^2 + L_1^2 - L_2^2}{2 L_1 D}\right)
\]
O resultado desta função é o vetor de ângulos alvo: $\mathbf{q}_d = \begin{bmatrix} q_{1d} \\ q_{2d} \end{bmatrix}$.

\section{Etapa 2: Geração de Trajetória por Interpolação}
Uma vez que $\mathbf{q}_d$ é conhecido, não o aplicamos instantaneamente, pois isso faria o robô "teleportar". Em vez disso, geramos uma trajetória suave no espaço de juntas usando \textbf{Interpolação Linear}.

Definimos um tempo de movimento $T$ (\texttt{MOVE\_DURATION}). O estado de animação $\tau$ é normalizado de 0 a 1:
\[
\tau(t) = \min\left( \frac{t_{\text{decorrido}}}{T}, 1.0 \right)
\]
A posição da junta interpolada $\mathbf{q}(t)$ é então calculada como:
\begin{equation}
\mathbf{q}(t) = \mathbf{q}_{\text{início}} + (\mathbf{q}_{\text{alvo}} - \mathbf{q}_{\text{início}}) \cdot \tau(t)
\label{eq:lerp}
\end{equation}
Isso corresponde às linhas de código: \texttt{t = trajectory\_time / MOVE\_DURATION} e \texttt{q\_current = q\_start + diff * t}.

\subsection{Tratamento de "Wrap-around" Angular}
Um problema crítico na interpolação de ângulos é o "wrap-around" de $2\pi$. Se $\mathbf{q}_{\text{início}} = -170^\circ$ e $\mathbf{q}_{\text{alvo}} = +170^\circ$, uma interpolação linear ingênua (Equação \ref{eq:lerp}) moveria o robô $340^\circ$ (o caminho longo).

Para garantir o caminho mais curto, a diferença de ângulo $\Delta \mathbf{q} = \mathbf{q}_{\text{alvo}} - \mathbf{q}_{\text{início}}$ deve ser normalizada para o intervalo $[-\pi, +\pi]$. Isso é feito no código:
\begin{itemize}
    \item \texttt{diff = q\_target - q\_start}
    \item \texttt{diff[i] = (diff[i] + np.pi) \% (2 * np.pi) - np.pi}
\end{itemize}
Esta nova \texttt{diff} representa o caminho angular mais curto, que é então usado na interpolação.

\section{Conclusão}
Esta abordagem em nível de posição é computacionalmente muito eficiente, pois a solução de CI analítica é calculada apenas uma vez. Ela não sofre de instabilidades de singularidade da mesma forma que os métodos baseados em Jacobiana.

A desvantagem é que esta é uma forma de controle em \textbf{malha aberta}. O robô simplesmente "executa" uma trajetória pré-calculada e não é reativo a perturbações ou a alvos que se movem dinamicamente.

\end{document}